%=====================================
%   20250528.tex
%   20250528 
%   toshi2019
%=====================================
% ------------------------
% documentclass
% ------------------------
\documentclass[leqno, 10pt, a4paper, dvipdfmx]{jsarticle}
\title{ゼミノート（06月25日分）}
\author{Toshi2019}
\date{\today}

% ------------------------
% usepackage
% ------------------------
\usepackage{algorithm}
\usepackage{algorithmic}
\usepackage{amscd}
\usepackage{amsfonts}
\usepackage{amsmath}
\usepackage[psamsfonts]{amssymb}
\usepackage{amsthm}
\usepackage{ascmac}
\usepackage{bm}
\usepackage{color}
\usepackage{enumerate}
\usepackage{fancybox}
\usepackage[stable]{footmisc}
\usepackage{graphicx}
\usepackage{listings}
\usepackage{mathrsfs}
\usepackage{mathtools}
\usepackage{otf}
\usepackage{pifont}
\usepackage{proof}
\usepackage{subfigure}
\usepackage{tikz}
\usepackage{verbatim}
\usepackage[all]{xy}

\usetikzlibrary{cd}
\usetikzlibrary{arrows.meta}


% ================================
% パッケージを追加する場合のスペース 
\usepackage[dvipdfmx]{hyperref}
\usepackage{xcolor}
\definecolor{darkgreen}{rgb}{0,0.45,0} 
\definecolor{darkred}{rgb}{0.75,0,0}
\definecolor{darkblue}{rgb}{0,0,0.6} 
\hypersetup{
    colorlinks=true,
    citecolor=darkgreen,
    linkcolor=darkred,
    urlcolor=darkblue,
}
\usepackage{pxjahyper}

%\usepackage[mathscr]{euscript} % mathscr のフォントを変更

%\usepackage{mmasym}

%=================================


% --------------------------
% theoremstyle
% --------------------------
\theoremstyle{definition}

% --------------------------
% newtheoem
% --------------------------

% 日本語で定理, 命題, 証明などを番号付きで用いるためのコマンドです. 
% If you want to use theorem environment in Japanece, 
% you can use these code. 
% Attention!
% All theorem enivironment numbers depend on 
% only section numbers.
\newtheorem{Axiom}{公理}[section]
\newtheorem{Definition}[Axiom]{定義}
\newtheorem{Theorem}[Axiom]{定理}
\newtheorem{Proposition}[Axiom]{命題}
\newtheorem{Lemma}[Axiom]{補題}
\newtheorem{Corollary}[Axiom]{系}
\newtheorem{Example}[Axiom]{例}
\newtheorem{Claim}[Axiom]{主張}
\newtheorem{Property}[Axiom]{性質}
\newtheorem{Attention}[Axiom]{注意}
\newtheorem{Question}[Axiom]{問}
\newtheorem{Problem}[Axiom]{問題}
\newtheorem{Consideration}[Axiom]{考察}
\newtheorem{Alert}[Axiom]{警告}
\newtheorem{Fact}[Axiom]{事実}


% 日本語で定理, 命題, 証明などを番号なしで用いるためのコマンドです. 
% If you want to use theorem environment with no number in Japanese, You can use these code.
\newtheorem*{Axiom*}{公理}
\newtheorem*{Definition*}{定義}
\newtheorem*{Theorem*}{定理}
\newtheorem*{Proposition*}{命題}
\newtheorem*{Lemma*}{補題}
\newtheorem*{Example*}{例}
\newtheorem*{Corollary*}{系}
\newtheorem*{Claim*}{主張}
\newtheorem*{Property*}{性質}
\newtheorem*{Attention*}{注意}
\newtheorem*{Question*}{問}
\newtheorem*{Problem*}{問題}
\newtheorem*{Consideration*}{考察}
\newtheorem*{Alert*}{警告}
\newtheorem{Fact*}{事実}


% 英語で定理, 命題, 証明などを番号付きで用いるためのコマンドです. 
%　If you want to use theorem environment in English, You can use these code.
%all theorem enivironment number depend on only section number.
\newtheorem{Axiom+}{Axiom}[section]
\newtheorem{Definition+}[Axiom+]{Definition}
\newtheorem{Theorem+}[Axiom+]{Theorem}
\newtheorem{Proposition+}[Axiom+]{Proposition}
\newtheorem{Lemma+}[Axiom+]{Lemma}
\newtheorem{Example+}[Axiom+]{Example}
\newtheorem{Corollary+}[Axiom+]{Corollary}
\newtheorem{Claim+}[Axiom+]{Claim}
\newtheorem{Property+}[Axiom+]{Property}
\newtheorem{Attention+}[Axiom+]{Attention}
\newtheorem{Question+}[Axiom+]{Question}
\newtheorem{Problem+}[Axiom+]{Problem}
\newtheorem{Consideration+}[Axiom+]{Consideration}
\newtheorem{Alert+}{Alert}
\newtheorem{Fact+}[Axiom+]{Fact}
\newtheorem{Remark+}[Axiom+]{Remark}

% ----------------------------
% commmand
% ----------------------------
% 執筆に便利なコマンド集です. 
% コマンドを追加する場合は下のスペースへ. 

% 集合の記号 (黒板文字)
\newcommand{\NN}{\mathbb{N}}
\newcommand{\ZZ}{\mathbb{Z}}
\newcommand{\QQ}{\mathbb{Q}}
\newcommand{\RR}{\mathbb{R}}
\newcommand{\CC}{\mathbb{C}}
\newcommand{\PP}{\mathbb{P}}
\newcommand{\KK}{\mathbb{K}}


% 集合の記号 (太文字)
\newcommand{\nn}{\mathbf{N}}
\newcommand{\zz}{\mathbf{Z}}
\newcommand{\qq}{\mathbf{Q}}
\newcommand{\rr}{\mathbf{R}}
\newcommand{\cc}{\mathbf{C}}
\newcommand{\pp}{\mathbf{P}}
\newcommand{\kk}{\mathbf{k}}

% 特殊な写像の記号
\newcommand{\ev}{\mathop{\mathrm{ev}}\nolimits} % 値写像
\newcommand{\pr}{\mathop{\mathrm{pr}}\nolimits} % 射影

% スクリプト体にするコマンド
%   例えば　{\mcal C} のように用いる
\newcommand{\mcal}{\mathcal}

% 花文字にするコマンド 
%   例えば　{\h C} のように用いる
\newcommand{\h}{\mathscr}

% ヒルベルト空間などの記号
\newcommand{\F}{\mcal{F}}
\newcommand{\X}{\mcal{X}}
\newcommand{\Y}{\mcal{Y}}
\newcommand{\Hil}{\mcal{H}}
\newcommand{\RKHS}{\Hil_{k}}
\newcommand{\Loss}{\mcal{L}_{D}}
\newcommand{\MLsp}{(\X, \Y, D, \Hil, \Loss)}

% 偏微分作用素の記号
\newcommand{\p}{\partial}

% 角カッコの記号 (内積は下にマクロがあります)
\newcommand{\lan}{\langle}
\newcommand{\ran}{\rangle}



% 圏の記号など
\newcommand{\Set}{{\bf Set}}
\newcommand{\Vect}{{\bf Vect}}
\newcommand{\FDVect}{{\bf FDVect}}
%\newcommand{\Ring}{{\bf Ring}}
\newcommand{\Ab}{{\bf Ab}}
\newcommand{\Mod}{\mathop{\mathrm{Mod}}\nolimits}
\newcommand{\Modf}{\mathop{\mathrm{Mod}^\mathrm{f}}\nolimits}
\newcommand{\CGA}{{\bf CGA}}
\newcommand{\GVect}{{\bf GVect}}
\newcommand{\Lie}{{\bf Lie}}
\newcommand{\dLie}{{\bf Liec}}



% 射の集合など
\newcommand{\Map}{\mathop{\mathrm{Map}}\nolimits} % 写像の集合
\newcommand{\Hom}{\mathop{\mathrm{Hom}}\nolimits} % 射集合
\newcommand{\End}{\mathop{\mathrm{End}}\nolimits} % 自己準同型の集合
\newcommand{\Aut}{\mathop{\mathrm{Aut}}\nolimits} % 自己同型の集合
\newcommand{\Mor}{\mathop{\mathrm{Mor}}\nolimits} % 射集合
\newcommand{\Ker}{\mathop{\mathrm{Ker}}\nolimits} % 核
\newcommand{\Img}{\mathop{\mathrm{Im}}\nolimits} % 像
\newcommand{\Cok}{\mathop{\mathrm{Coker}}\nolimits} % 余核
\newcommand{\Cim}{\mathop{\mathrm{Coim}}\nolimits} % 余像

% その他便利なコマンド
\newcommand{\dip}{\displaystyle} % 本文中で数式モード
\newcommand{\e}{\varepsilon} % イプシロン
\newcommand{\dl}{\delta} % デルタ
\newcommand{\pphi}{\varphi} % ファイ
\newcommand{\ti}{\tilde} % チルダ
\newcommand{\pal}{\parallel} % 平行
\newcommand{\op}{{\rm op}} % 双対を取る記号
\newcommand{\lcm}{\mathop{\mathrm{lcm}}\nolimits} % 最小公倍数の記号
\newcommand{\Probsp}{(\Omega, \F, \P)} 
\newcommand{\argmax}{\mathop{\rm arg~max}\limits}
\newcommand{\argmin}{\mathop{\rm arg~min}\limits}





% ================================
% コマンドを追加する場合のスペース 
%\newcommand{\OO}{\mcal{O}}



\makeatletter
\renewenvironment{proof}[1][\proofname]{\par
  \pushQED{\qed}%
  \normalfont \topsep6\p@\@plus6\p@\relax
  \trivlist
  \item[\hskip\labelsep
%        \itshape
         \bfseries
%    #1\@addpunct{.}]\ignorespaces
    {#1}]\ignorespaces
}{%
  \popQED\endtrivlist\@endpefalse
}
\makeatother

\renewcommand{\proofname}{証明.}



%\renewcommand\proofname{\bf 証明} % 証明
\numberwithin{equation}{section}
\newcommand{\cTop}{\textsf{Top}}
%\newcommand{\cOpen}{\textsf{Open}}
\newcommand{\Op}{\mathop{\textsf{Op}}\nolimits}
\newcommand{\Ob}{\mathop{\textrm{Ob}}\nolimits}
\newcommand{\id}{\mathop{\mathrm{id}}\nolimits}
\newcommand{\pt}{\mathop{\mathrm{pt}}\nolimits}
\newcommand{\res}{\mathop{\rho}\nolimits}
\newcommand{\A}{\mcal{A}}
\newcommand{\B}{\mcal{B}}
\newcommand{\C}{\mcal{C}}
\newcommand{\D}{\mcal{D}}
\newcommand{\E}{\mcal{E}}
\newcommand{\G}{\mcal{G}}
%\newcommand{\H}{\mcal{H}}
\newcommand{\I}{\mcal{I}}
\newcommand{\J}{\mcal{J}}
\newcommand{\OO}{\mcal{O}}
\newcommand{\Ring}{\mathop{\textsf{Ring}}\nolimits}
\newcommand{\cAb}{\mathop{\textsf{Ab}}\nolimits}
%\newcommand{\Ker}{\mathop{\mathrm{Ker}}\nolimits}
\newcommand{\im}{\mathop{\mathrm{Im}}\nolimits}
\newcommand{\Coker}{\mathop{\mathrm{Coker}}\nolimits}
\newcommand{\Coim}{\mathop{\mathrm{Coim}}\nolimits}
\newcommand{\rank}{\mathop{\mathrm{rank}}\nolimits}
\newcommand{\Ht}{\mathop{\mathrm{Ht}}\nolimits}
\newcommand{\supp}{\mathop{\mathrm{supp}}\nolimits}
\newcommand{\colim}{\mathop{\mathrm{colim}}}
\newcommand{\Tor}{\mathop{\mathrm{Tor}}\nolimits}

\newcommand{\cat}{\mathscr{C}}

%筆記体
\newcommand{\cA}{\mcal{A}}
\newcommand{\cB}{\mcal{B}}
\newcommand{\cC}{\mcal{C}}
\newcommand{\cD}{\mcal{D}}
\newcommand{\cE}{\mcal{E}}
\newcommand{\cF}{\mcal{F}}
\newcommand{\cG}{\mcal{G}}
\newcommand{\cH}{\mcal{H}}
\newcommand{\cI}{\mcal{I}}
\newcommand{\cJ}{\mcal{J}}
\newcommand{\cK}{\mcal{K}}
\newcommand{\cL}{\mcal{L}}
\newcommand{\cM}{\mcal{M}}
\newcommand{\cN}{\mcal{N}}
\newcommand{\cO}{\mcal{O}}
\newcommand{\cP}{\mcal{P}}
\newcommand{\cQ}{\mcal{Q}}
\newcommand{\cR}{\mcal{R}}
\newcommand{\cS}{\mcal{S}}
\newcommand{\cT}{\mcal{T}}
\newcommand{\cU}{\mcal{U}}
\newcommand{\cV}{\mcal{V}}
\newcommand{\cW}{\mcal{W}}
\newcommand{\cX}{\mcal{X}}
\newcommand{\cY}{\mcal{Y}}
\newcommand{\cZ}{\mcal{Z}}


\newcommand{\scA}{\mathscr{A}}
\newcommand{\scB}{\mathscr{B}}
\newcommand{\scC}{\mathscr{C}}
\newcommand{\scD}{\mathscr{D}}
\newcommand{\scE}{\mathscr{E}}
\newcommand{\scF}{\mathscr{F}}
\newcommand{\scN}{\mathscr{N}}
\newcommand{\scO}{\mathscr{O}}
\newcommand{\scV}{\mathscr{V}}
\newcommand{\scU}{\mathscr{U}}


\newcommand{\ibA}{\mathop{\text{\textit{\textbf{A}}}}}
\newcommand{\ibB}{\mathop{\text{\textit{\textbf{B}}}}}
\newcommand{\ibC}{\mathop{\text{\textit{\textbf{C}}}}}
\newcommand{\ibD}{\mathop{\text{\textit{\textbf{D}}}}}
\newcommand{\ibE}{\mathop{\text{\textit{\textbf{E}}}}}
\newcommand{\ibF}{\mathop{\text{\textit{\textbf{F}}}}}
\newcommand{\ibG}{\mathop{\text{\textit{\textbf{G}}}}}
\newcommand{\ibH}{\mathop{\text{\textit{\textbf{H}}}}}
\newcommand{\ibI}{\mathop{\text{\textit{\textbf{I}}}}}
\newcommand{\ibJ}{\mathop{\text{\textit{\textbf{J}}}}}
\newcommand{\ibK}{\mathop{\text{\textit{\textbf{K}}}}}
\newcommand{\ibL}{\mathop{\text{\textit{\textbf{L}}}}}
\newcommand{\ibM}{\mathop{\text{\textit{\textbf{M}}}}}
\newcommand{\ibN}{\mathop{\text{\textit{\textbf{N}}}}}
\newcommand{\ibO}{\mathop{\text{\textit{\textbf{O}}}}}
\newcommand{\ibP}{\mathop{\text{\textit{\textbf{P}}}}}
\newcommand{\ibQ}{\mathop{\text{\textit{\textbf{Q}}}}}
\newcommand{\ibR}{\mathop{\text{\textit{\textbf{R}}}}}
\newcommand{\ibS}{\mathop{\text{\textit{\textbf{S}}}}}
\newcommand{\ibT}{\mathop{\text{\textit{\textbf{T}}}}}
\newcommand{\ibU}{\mathop{\text{\textit{\textbf{U}}}}}
\newcommand{\ibV}{\mathop{\text{\textit{\textbf{V}}}}}
\newcommand{\ibW}{\mathop{\text{\textit{\textbf{W}}}}}
\newcommand{\ibX}{\mathop{\text{\textit{\textbf{X}}}}}
\newcommand{\ibY}{\mathop{\text{\textit{\textbf{Y}}}}}
\newcommand{\ibZ}{\mathop{\text{\textit{\textbf{Z}}}}}

\newcommand{\ibx}{\mathop{\text{\textit{\textbf{x}}}}}

%\newcommand{\Comp}{\mathop{\mathrm{C}}\nolimits}
%\newcommand{\Komp}{\mathop{\mathrm{K}}\nolimits}
%\newcommand{\Domp}{\mathop{\mathsf{D}}\nolimits}%複体のホモトピー圏
%\newcommand{\Comp}{\mathrm{C}}
%\newcommand{\Komp}{\mathrm{K}}
%\newcommand{\Domp}{\mathsf{D}}%複体のホモトピー圏

\newcommand{\Comp}{\mathop{\mathrm{C}}\nolimits}
\newcommand{\Komp}{\mathop{\mathsf{K}}\nolimits}
\newcommand{\Domp}{\mathord{\mathsf{D}}}%\nolimits}
\newcommand{\Kompl}{\mathop{\mathsf{K}^\mathrm{+}}\nolimits}
\newcommand{\Kompu}{\mathop{\mathsf{K}^\mathrm{-}}\nolimits}
\newcommand{\Kompb}{\mathop{\mathsf{K}^\mathrm{b}}\nolimits}
\newcommand{\Dompl}{\mathop{\mathsf{D}^\mathrm{+}}\nolimits}
\newcommand{\Dompu}{\mathop{\mathsf{D}^\mathrm{-}}\nolimits}
\newcommand{\Dompb}{\mathop{\mathsf{D}^\mathrm{b}}\nolimits}
\newcommand{\Dbconic}{\mathop{\mathsf{D}^\mathrm{b}_{\rrp}}\nolimits}
\newcommand{\Dompbf}{\mathop{\mathsf{D}_\mathrm{f}^\mathrm{b}}\nolimits}
\newcommand{\Domplb}{\mathop{\mathsf{D}^\mathrm{lb}}\nolimits}




\newcommand{\CCat}{\Comp(\cat)}
\newcommand{\KCat}{\Komp(\cat)}
\newcommand{\DCat}{\Domp(\cat)}%圏Cの複体のホモトピー圏
%\newcommand{\HOM}{\mathop{\mathscr{H}\hspace{-2pt}om}\nolimits}%内部Hom
\newcommand{\HOM}{\mathop{\mathcal{H}\hspace{-0pt}om}\nolimits}%内部Hom
\newcommand{\RHOM}{\mathop{\mathrm{R}\hspace{-1.5pt}\HOM}\nolimits}

\newcommand{\muS}{\mathop{\mathrm{SS}}\nolimits}
\newcommand{\RG}{\mathop{\mathrm{R}\hspace{-0pt}\Gamma}\nolimits}
\newcommand{\RHom}{\mathop{\mathrm{R}\hspace{-1.5pt}\Hom}\nolimits}
\newcommand{\Rder}{\mathrm{R}}

\newcommand{\simar}{\mathrel{\overset{\sim}{\rightarrow}}}%同型右矢印
\newcommand{\simarr}{\mathrel{\overset{\sim}{\longrightarrow}}}%同型右矢印
\newcommand{\simra}{\mathrel{\overset{\sim}{\leftarrow}}}%同型左矢印
\newcommand{\simrra}{\mathrel{\overset{\sim}{\longleftarrow}}}%同型左矢印

\newcommand{\hocolim}{{\mathrm{hocolim}}}
\newcommand{\indlim}[1][]{\mathop{\varinjlim}\limits_{#1}}
\newcommand{\sindlim}[1][]{\smash{\mathop{\varinjlim}\limits_{#1}}\,}
\newcommand{\Pro}{\mathrm{Pro}}
\newcommand{\Ind}{\mathrm{Ind}}
\newcommand{\prolim}[1][]{\mathop{\varprojlim}\limits_{#1}}
\newcommand{\sprolim}[1][]{\smash{\mathop{\varprojlim}\limits_{#1}}\,}

\newcommand{\Sh}{\mathrm{Sh}}
\newcommand{\PSh}{\mathrm{PSh}}

\newcommand{\rmD}{\mathsf{D}}
\newcommand{\vD}{\mathbf{D}}

\newcommand{\Lloc}[1][]{\mathord{\mathcal{L}^1_{\mathrm{loc},{#1}}}}
\newcommand{\ori}{\mathord{\mathrm{or}}}
\newcommand{\Db}{\mathord{\mathscr{D}b}}

\newcommand{\codim}{\mathop{\mathrm{codim}}\nolimits}



\newcommand{\gld}{\mathop{\mathrm{gld}}\nolimits}
\newcommand{\wgld}{\mathop{\mathrm{wgld}}\nolimits}


\newcommand{\tens}[1][]{\mathbin{\otimes_{\raise1.5ex\hbox to-.1em{}{#1}}}}
\newcommand{\etens}{\mathbin{\boxtimes}}
\newcommand{\ltens}[1][]{\mathbin{\overset{\mathrm{L}}\tens}_{#1}}
\newcommand{\mtens}[1][]{\mathbin{\overset{\mathrm{\mu}}\tens}_{#1}}
\newcommand{\lltens}[1][]{\mathbin{{\mathop{\tens}\limits^{\mathrm{L}}_{#1}}}}
\newcommand{\lletens}[1][]{\mathbin{{\mathop{\boxtimes}\limits^{\mathrm{L}}_{#1}}}}
\newcommand{\letens}{\overset{\mathrm{L}}{\etens}}
\newcommand{\detens}{\underline{\etens}}
\newcommand{\ldetens}{\overset{\mathrm{L}}{\underline{\etens}}}
\newcommand{\dtens}[1][]{{\overset{\mathrm{L}}{\underline{\otimes}}}_{#1}}

\newcommand{\blk}{\mathord{\ \cdot\ }}
\newcommand{\mres}[2][]{{\left.{#1}\right\rvert}_{#2}}


%\newcommand{\hocolim}{{\mathrm{hocolim}}}
%\newcommand{\indlim}[1][]{\mathop{\varinjlim}\limits_{#1}}
%\newcommand{\sindlim}[1][]{\smash{\mathop{\varinjlim}\limits_{#1}}\,}
%\newcommand{\Pro}{\mathrm{Pro}}
%\newcommand{\Ind}{\mathrm{Ind}}
%\newcommand{\prolim}[1][]{\mathop{\varprojlim}\limits_{#1}}
%\newcommand{\sprolim}[1][]{\smash{\mathop{\varprojlim}\limits_{#1}}\,}
\newcommand{\proolim}[1][]{\mathop{\text{\rm``{$\varprojlim$}''}}\limits_{#1}}
\newcommand{\sproolim}[1][]{\smash{\mathop{\rm``{\varprojlim}''}\limits_{#1}}}
\newcommand{\inddlim}[1][]{\mathop{\text{\rm``{$\varinjlim$}''}}\limits_{#1}}
\newcommand{\sinddlim}[1][]{\smash{\mathop{\text{\rm``{$\varinjlim$}''}}\limits_{#1}}\,}
\newcommand{\ooplus}{\mathop{\text{\rm``{$\oplus$}''}}\limits}
\newcommand{\bbigsqcup}{\mathop{``\bigsqcup''}\limits}
\newcommand{\bsqcup}{\mathop{``\sqcup''}\limits}
\newcommand{\dsum}[1][]{\mathbin{\oplus_{#1}}}

\newcommand{\Fct}{\mathop{\mathsf{Fct}}\nolimits}

\newcommand{\pb}[1][]{\mathbin{\mathop{\times}\limits_{#1}}}
\newcommand{\dpb}[1][]{\mathbin{\mathop{\times}\limits_{#1}}}
\newcommand{\tpb}[1][]{\mathbin{\mathop{\times}\nolimits_{#1}}}





%================================================
% 自前の定理環境
%   https://mathlandscape.com/latex-amsthm/
% を参考にした
\newtheoremstyle{mystyle}%   % スタイル名
    {5pt}%                   % 上部スペース
    {5pt}%                   % 下部スペース
    {}%              % 本文フォント
    {}%                  % 1行目のインデント量
    {\bfseries}%                      % 見出しフォント
    {.}%                     % 見出し後の句読点
    {12pt}%                     % 見出し後のスペース
    {\thmname{#1}\thmnumber{ #2}\thmnote{{\hspace{2pt}\normalfont (#3)}}}% % 見出しの書式

\theoremstyle{mystyle}
\newtheorem{AXM}{公理}[section]
\newtheorem{DFN}[AXM]{定義}
\newtheorem{THM}[AXM]{定理}
\newtheorem*{THM*}{定理}
\newtheorem{PRP}[AXM]{命題}
\newtheorem{LMM}[AXM]{補題}
\newtheorem{CRL}[AXM]{系}
\newtheorem{EG}[AXM]{例}
\newtheorem*{EG*}{例}
\newtheorem{RMK}[AXM]{注意}
\newtheorem{CNV}[AXM]{約束}
\newtheorem{CMT}[AXM]{コメント}
\newtheorem*{CMT*}{コメント}
\newtheorem{NTN}[AXM]{記号}

% 定理環境ここまで
%====================================================

\usepackage{framed}
\definecolor{lightgray}{rgb}{0.75,0.75,0.75}
\renewenvironment{leftbar}{%
  \def\FrameCommand{\textcolor{lightgray}{\vrule width 4pt} \hspace{10pt}}% 
  \MakeFramed {\advance\hsize-\width \FrameRestore}}%
{\endMakeFramed}
\newenvironment{redleftbar}{%
  \def\FrameCommand{\textcolor{lightgray}{\vrule width 1pt} \hspace{10pt}}% 
  \MakeFramed {\advance\hsize-\width \FrameRestore}}%
 {\endMakeFramed}



\renewcommand{\Re}{\mathop{\mathrm{Re}}\nolimits}

% =================================





% ---------------------------
% new definition macro
% ---------------------------
% 便利なマクロ集です

% 内積のマクロ
%   例えば \inner<\pphi | \psi> のように用いる
\def\inner<#1>{\langle #1 \rangle}

% ================================
% マクロを追加する場合のスペース 

%=================================



% 位相空間
\newcommand{\Int}[1][]{\mathop{\mathrm{Int}}\nolimits_{#1}}
\newcommand{\Cl}[1][]{\mathop{\mathrm{Cl}}\nolimits_{#1}}
\newcommand{\Bd}[1][]{\mathop{\mathrm{Bd}}\nolimits_{#1}}
\newcommand{\bd}{\mathord{\partial}}
\newcommand{\Fr}[1][]{\mathop{\mathrm{Fr}}\nolimits_{#1}}


\newcommand{\mtan}[1][]{T\hspace{-1.75pt}{#1}}
\newcommand{\mtp}[1][]{{}^{t\!}{#1}} % transpose of the morphism
%\newcommand{\mpb}[1][]{{}^{t\!}{#1}'} % pull-back of the morphism
\newcommand{\mpb}[1][]{{#1}_{\pi}} % pull-back of the morphism
\newcommand{\mdiff}[1][]{{#1}_{d}} % pull-back of the morphism
\newcommand{\mptan}[1][]{{#1}_{\tau}} % pull-back of the morphism


%\newcommand{\pb}[1][]{\mathbin{\mathop{\times}\limits_{#1}}}

\newcommand{\cotb}[1][]{T^{\ast}{#1}}
\newcommand{\conm}[2][]{T_{#1}^{\hspace{0.7pt}\ast}{#2}}
\newcommand{\norb}[2][]{T_{#1}{#2}}

\newcommand{\cotz}[1][]{T^{\ast}_{#1}{#1}}

\newcommand{\cotn}[1][]{\dot{T}^{\ast}{#1}}


\newcommand{\muhom}{\mu hom}
\newcommand{\muHom}[1][]{\mathrm{Hom}^\mu_{\raise1.5ex\hbox to.1em{}#1}}

\newcommand{\mucat}[2][]{\mathsf{D}^{\mathrm{b}}_{#1}(#2)}
%\newcommand{\mucat}[2][]{\mathsf{D}^{\mathrm{b}}_{#1}(#2)}

\newcommand{\conv}[1][]{\mathbin{\mathop{\circ}\limits_{#1}}}

\newcommand{\torus}{\mathbf{T}}

\newcommand{\diag}[1][]{\Delta_{#1}}
%  \[\mu_{\Delta_{Y}}\mathrm{R}\mathcal{H}om(G,F)\]
\newcommand{\micro}[1]{\mu_{#1}\hspace{-0pt}}

% ----------------------------
% documenet 
% ----------------------------
% 以下, 本文の執筆スペースです. 
% Your main code must be written between 
% begin document and end document.
% ---------------------------

\begin{document}
\maketitle
\tableofcontents

\section{復習}

\subsection*{双対化複体のコホモロジーが集中していること}

\(X\)を多様体，\(Y\)を余次元\(n\)の閉部分多様体とし，
\(j\colon Y\hookrightarrow X\)を閉うめこみとする．
このとき，\(\omega_{Y/X}=j^{!}A_X\)のコホモロジーを計算する．
まず，\(j^{!}A_X=j^{-1}\RG_{Y}A_X\)である．

\(j^{-1}\cO_X\tens[]\omega_{M/X}\to j^{!}\cO_X\)

\(j^{-1}A_X\tens[]j^{!}A_X\to j^{!}A_X\)

\(j^{-1}\zz_X\tens[]j^{!}\zz_X\to j^{!}\zz_X\)

\(j^{-1}\zz_X\tens[]j^{-1}\RG_{M}\zz_X\to j^{-1}\RG_{M}\zz_X\)
\(\RG_{M}\zz_X\)の\(x\in M\)での茎を見ればよい．
完全三角\[
    \RG_{M}\zz_{X}\to \zz_{X}\to \RG_{X-M}\zz_{X}\to +1
\]から，長完全列
\begin{align*}
    \RG_{M}\zz_{X}&\to \zz_{X}\to \RG_{X-M}\zz_{X}\\
    \to H^{1}_{M}(\zz_X)&\to H^{1}(\zz_{X})\to H^{1}_{X-M}(\zz_X)\\
    \to\cdots
\end{align*}


\section{関手\(\muhom\)}
%\newcommand{\mucat}[2][]{\mathsf{D}^{\mathrm{b}}_{#1}(#2)}
\(f\)を\(Y\)から\(X\)への射とする．
\(\diag[f]\subset X\times Y\)を\(f\)のグラフとする．
\(\diag[X]\)と\(\diag[Y]\)で
それぞれ\(X\times X\)と
\(Y\times Y\)の対角集合を表す．
\(Y\mathbin{\times_{X}}\cotb[X]\)と\(T^{\ast}_{\diag[f]}(X\times Y)\)を
射影\(\cotb[(X\times Y)]\to(\cotb[X])\times Y\)で同一視する．
同様に\(\cotb[X]\)と\(T^{\ast}_{\diag[X]}(X\times X)\)，
\(\cotb[Y]\)と\(T^{\ast}_{\diag[Y]}(Y\times Y)\)を
それぞれ第1射影で同一視する．
混同の恐れがないときは\(\diag[f]\)を\(\diag\)とかく．
すると，次の写像を得る（式 (4.3.3) 参照）．
\begin{equation}\label{eq:4.4.1}
  \vcenter{\xymatrix@C=36pt@R=36pt{
    T^{\ast}_{\diag[Y]}(Y\times Y)
    \ar[d]^-{\rotatebox{90}{\(\sim\)}}
    &
    T^{\ast}_{\diag[f]}(X\times Y)
    \ar[d]^-{\rotatebox{90}{\(\sim\)}}
    \ar[l]
    \ar[r]
    &
    T^{\ast}_{\diag[X]}(X\times X)
    \ar[d]^-{\rotatebox{90}{\(\sim\)}}
    \\
    \cotb[Y]
    &
    Y\pb[X]\cotb[X]
    \ar[l]^-{\mdiff[f]}
    \ar[r]_-{\mpb[f]}
    &
    \cotb[X].
  }}
\end{equation}
次の公式をよく用いる．
\begin{equation}\label{eq:4.4.2}
  \omega_{{Y\mathbin{\times_{X}}\cotb[X]}\mathbin{/}\cotb[Y]}
  \tens[]
  \omega_{{Y\mathbin{\times_{X}}\cotb[X]}\mathbin{/}\cotb[X]}
  \cong
  A_{Y\mathbin{\times_{X}}\cotb[X]}
\end{equation}
\begin{proof}
  \(
    \omega_{{Y\mathbin{\times_{X}}\cotb[X]}\mathbin{/}\cotb[Y]}
    \tens[]
    \omega_{{Y\mathbin{\times_{X}}\cotb[X]}\mathbin{/}\cotb[X]}
  \)は\(A_{Y\mathbin{\times_{X}}\cotb[X]}\)加群の複体なので射
  \[
    A_{Y\mathbin{\times_{X}}\cotb[X]}
    \to
    \omega_{{Y\mathbin{\times_{X}}\cotb[X]}\mathbin{/}\cotb[Y]}
    \tens[]
    \omega_{{Y\mathbin{\times_{X}}\cotb[X]}\mathbin{/}\cotb[X]}
  \]
  が存在する．したがって茎ごとに同型であることを示せばよい．
  \((y,f(x);\xi)\)を\(Y\mathbin{\times_{X}}\cotb[X]\)の点とする．
  \begin{align*}
    \left(
      \mdiff[f]^{!}A_{\cotb[Y]}\tens[]\mpb[f]^{!}A_{\cotb[X]}
    \right)_{(y,f(x);\xi)}
    &=\left(
      \mdiff[f]^{-1}A_{\cotb[Y]}\tens[]\mpb[f]^{-1}A_{\cotb[X]}
    \right)_{(y,f(x);\xi)}\\
    &=\left(\mdiff[f]^{-1}A_{\cotb[Y]}\right)_{(y,f(x);\xi)}
    \tens[]\left(\mpb[f]^{-1}A_{\cotb[X]}
    \right)_{(y,f(x);\xi)}\\
    &=\left(A_{\cotb[Y]}\right)_{(y;f^{\ast}\xi)}
    \tens[]\left(\mpb[f]^{-1}A_{\cotb[X]}
    \right)_{(y;f^{\ast}\xi)}\\
    &=A\\
    &=\left(A_{Y\mathbin{\times_{X}}\cotb[X]}\right)_{(y,f(x);\xi)}
  \end{align*}
  である．
\end{proof}

\(f_1\colon Y\times Y\to X\times Y\)で写像\(f\times \id_Y\)を表し，
\(f_1\colon X\times Y\to X\times X\)で写像\(\id_X\times f\)を表す．
これにより以下の可換図式を得る．
ただし右側の四角はカルテシアンであり\(f_2\)は\(\diag[X]\)に対して横断的である．
\begin{equation}\label{eq:4.4.3}
  \vcenter{\xymatrix@C=36pt@R=36pt{
    Y\times Y
    \ar[r]_-{f_1}
    &
    X\times Y
    \ar@{}[rd]|\square
    %\ar[d]^-{\rotatebox{90}{\(\sim\)}}
    \ar[r]_-{f_2}
    &
    X\times X
    %\ar[d]^-{\rotatebox{90}{\(\sim\)}}
    \\
    \diag[Y]  
    \ar@{^{(}->}[u]^-{}
    \ar[r]^-{{\sim}}
    &
    \diag[]
    \ar@{^{(}->}[u]^-{}
    %\ar[l]^-{\mdiff[f]}
    \ar[r]_-{f}
    &
    \diag[X].
    \ar@{^{(}->}[u]^-{}
  }}
\end{equation}
\begin{proof}
  カルテシアンであることは閉うめこみから従う．

  \(f_2\)が\(\diag[X]\)に対して横断的であることを示す．
  \(\mres[{\mdiff[f_2]}]{(X\times Y)\pb[X\times X]T^{\ast}_{\diag[X]}(X\times X)}\)が単射であればよい．
  %この写像は局所的に\[
  %  (f(y),y;\xi,0)\mapsto (f(y),y;\xi,0)
  %\]とかけるので完了．
  まず，\(\mdiff[f_2]\)は
  \[    
    \vcenter{\xymatrix@C=26pt@R=10pt{
      (X\times Y)\pb[X\times X]T^{\ast}(X\times X)
      %\ar[d]_-{\tilde{f'}}
      \ar[r]^-{\mdiff[f_2]}
      %\ar@{}[dr]|\square
      &
      T^{\ast}(X\times Y)
      %\ar[d]^-{f}
      \\
      \left(\left(f(y),y\right),(f(y),f(y);\xi_1,\xi_2)\right)
      \ar@{}[u]|{\rotatebox{90}{\(\in\)}}
      \ar@{{|}->}[r]_-{}
      %\ar@{}[dr]|\square
      &
      \ar@{}[u]|{\rotatebox{90}{\(\in\)}}
      \left(f(y),y;f_{2}^{\ast}(\xi_1,\xi_2)\right)
    }}
  \]であり，
  \[
    f_2^{\ast}(\xi_1,\xi_2)
    =
    (\id_X\times f)^{\ast}(\xi_1,\xi_2)
    =(\xi_1,f^{\ast}\xi_2)
  \]であるから，\(T^{\ast}_{\diag[X]}(X\times X)\)の点が局所的に
  \((x,x;\xi,-\xi)\)とかけることに注意すると，
  \[
    \mres[{\mdiff[f_2]}]{(X\times Y)\pb[X\times X]T^{\ast}_{\diag[X]}(X\times X)}
    \left(\left(f(y),y\right),(f(y),f(y);\xi,-\xi)\right)
    =(f(y),y;\xi,-f^{\ast}\xi)
  \]
  とかける．\((f(y),y;\xi,-\xi)\)に対し，\((f(y),y;\xi,-f^{\ast}\xi)=(f(y),y;0,0)\)となるとすると，
  \(\xi=0,-f^{\ast}\xi=0\)より，\(\xi=-\xi=0\)となるので単射である．
\end{proof}
次の記号を定める．
\begin{itemize}
  \item \(q_j\)：\(X\times X\)上の第\(j\)射影
  \item \(\tilde{q}_{j}\)：\(X\times Y\)上の第\(j\)射影
  \item \(q'_j\)：\(Y\times Y\)上の第\(j\)射影
\end{itemize}

\begin{DFN}\label{dfn4.4.1}
  \(G\in\Dompb(Y)\), \(F\in\Dompb(X)\)に対し，次のようにおく．
  \begin{enumerate}[(i)]
    \item \(
      \muhom(G\to F)
      \coloneqq
      \micro{\diag[]}\RHOM_{
        %A_{X\times Y}
        }(\tilde{q}_2^{-1}G,\tilde{q}_1^{!}F)\)
    \item \(
      \muhom(F\leftarrow G)
      \coloneqq
      \left(
        \micro{\diag[]}\RHOM_{
          %A_{X\times Y}
        }(\tilde{q}_2^{-1}G,\tilde{q}_1^{!}F)
      \right)^{a}\)
    \item \(Y\to X\)で\(f\)が恒等射のとき，\[
      \muhom(G,F)
      \coloneqq\muhom(G\to F)
      =\micro{\diag[X]}\RHOM_{
        %A_{X\times X}
      }({q}_2^{-1}G,{q}_1^{!}F)
    \]
  \end{enumerate}
\end{DFN}
(ii)の右辺の\((\blk)^{a}\)は\(Y\tpb[X]\cotb[X]\)上の
対蹠写像による逆像を表す．

\(\pi\)で\(T^{\ast}_{\diag[]}(X\times Y)\)から\(\diag[]\cong Y\)への射影を表す．

\begin{leftbar}
\begin{PRP}\label{prp4.4.2}
  次の標準的な同型が存在する．
  \begin{enumerate}[(i)]
    \item 一般に
    %\vspace{-2\baselineskip}
    \begin{align*}
      %\vspace{-2\baselineskip}
      \Rder{\pi_{\ast}}\muhom(G\to F)&\cong\RHOM(G,f^{!}F),\\
      %\vspace{-2\baselineskip}
      \Rder{\pi_{\ast}}\muhom(F\leftarrow G)&\cong\RHOM(f^{-1}F,G).
    \end{align*}
    \(Y=X\)で\(f\)が恒等射の場合はとくに\[
      \Rder{\pi_{\ast}}\muhom(G, F)\cong\RHOM(G,F)
    \]
    となる．
    \item \(G\)がコホモロジー構成可能層と仮定する．このとき次が成り立つ．\[
      \Rder{\pi_{!}}\muhom(G\to F)
      \cong
      \RHOM(G,A_Y)\lltens[]f^{-1}F\tens[]\omega_{Y/X}.
    \]
    \(F\)がコホモロジー構成可能層のときは
    \[
      \Rder{\pi_{!}}\muhom(F\leftarrow G)
      \cong
      f^{-1}\RHOM(f^{-1}F,A_X)\lltens[]G.
    \]が成り立つ．
    \(Y=X\)で\(f\)が恒等射の場合はとくに\[
      \Rder{\pi_{!}}\muhom(G, F)\cong\RHOM(G,A_X)\lltens[]F
    \]
    となる．
  \end{enumerate}
\end{PRP}
\end{leftbar}
\begin{proof}
  (i) 
  定理4.3.2から\begin{align*}
    \Rder{\pi_{\ast}}\micro{\diag[]}\RHOM(\tilde{q}_2^{-1}G,\tilde{q}_1^{!}F)
    &\underset{(1)}{\cong}
    \Rder{{\tilde{q_2}}_{\ast}}\RG_{\diag[]}\RHOM(\tilde{q}_2^{-1}G,\tilde{q}_1^{!}F)\\
    &\underset{(2)}{\cong}
    \Rder{{\tilde{q_2}}_{\ast}}\RHOM(\tilde{q}_2^{-1}G,\RG_{\diag[]}\tilde{q}_1^{!}F)\\
    &\underset{(3)}{\cong}
    \RHOM(G,\Rder{{\tilde{q_2}}_{\ast}}\RG_{\diag[]}\tilde{q}_1^{!}F)\\
    &\underset{(4)}{\cong}
    \RHOM(G,f^{!}F)
  \end{align*}
  である．ただし，それぞれの同型の根拠は次のとおり．
  \begin{enumerate}[(1)]
    \item 定理4.3.2 (iv) の同型\(
      \mres[\mu_M(F)]{M}\cong i^{!}F
    \)と同一視\(Y\cong\diag[]\)から従う．
    詳しくは以下のようになる．
    \(i\colon \diag[]\hookrightarrow X\times Y\)を閉うめこみ，
    \(\varphi\colon\diag[]\to Y\)を射影，
    \(\psi\colon Y\to \diag[]\)をその逆写像とする．
    このとき，\begin{align*}
      \Rder{\pi_{\ast}}\mu_{\diag[]}
      &\cong
      i^{!}&&\text{（ここに定理4.3.2）}\\
      &\cong
      \Rder{\varphi_{\ast}}i^{!}&&\text{（同一視からくる同型）}\\
      &\cong
      \Rder{\tilde{q_2}_{\ast}}\Rder{i_{\ast}}i^{!}&&\text{（図式の可換性）}\\
      &\cong
      \Rder{\tilde{q_2}_{\ast}}\RG_{\diag[]}&&\text{（\(\RG_{\diag[]}\)の書き換え）}
    \end{align*}
    である．
    \item \(\RG_{\diag[]}\)に関する公式(2.6.9)
    \item \(\tilde{q}_2\)に関する順像逆像随伴
    \item これも同一視からくる．詳しくかくと，
    図式\[\vcenter{\xymatrix{
      &
      X\times Y
      \ar[ld]_-{\tilde{q}_1}
      \ar[rd]^-{\tilde{q}_2}&\\
      X&
      \diag[]
      \ar[l]^-{f}
      \ar[u]_-{i}
      \ar[r]_-{\varphi}
      &Y
    }}\]を用いて，
    \begin{align*}
      f^{!}F
      &\cong i^{!}\tilde{q_1}^{!}F&&\text{（図式の可換性）}\\
      &\cong \Rder{\varphi_{\ast}}i^{!}\tilde{q_1}^{!}F&&\text{（同一視からくる同型）}\\
      &\cong \Rder{\tilde{q_2}_{\ast}}\Rder{i_{\ast}}i^{!}\tilde{q_1}^{!}F&&\text{（図式の可換性）}\\
      &\cong \Rder{\tilde{q_2}_{\ast}}\RG_{\diag[]}\tilde{q_1}^{!}F&&\text{（\(\RG_{\diag[]}\)の書き換え）}
    \end{align*}
    のようになる．
  \end{enumerate}
  2つ目の式は，\(\pi=\pi\circ a\)に注意すれば同様に
  \begin{align*}
    \Rder{\pi_{\ast}}(
      \micro{\diag[]}\RHOM(
        \tilde{q}_1^{-1}F,\tilde{q}_{2}^{!}G
      )
    )^{a}
    &\underset{}{\cong}
    \Rder{\pi_{\ast}}\micro{\diag[]}\RHOM(
      \tilde{q}_1^{-1}F,\tilde{q}_{2}^{!}G
    )\\
    &\underset{}{\cong}
    \Rder{{\tilde{q_2}}_{\ast}}\RG_{\diag}\RHOM(\tilde{q}_{1}^{-1}F,\tilde{q}_2^{!}G)\\
    &\underset{(2.6.9)}{\cong}
    \Rder{{\tilde{q_2}}_{\ast}}\RHOM((\tilde{q}_1^{-1}F)_{\diag[]},\tilde{q}_2^{!}G)\\
    &\underset{(V.D.)}{\cong}
    \RHOM(\Rder{{\tilde{q_2}}_{!}}(\tilde{q}_1^{-1}F)_{\diag[]},G)\\
    &\underset{}{\cong}
    \RHOM(f^{-1}F,G)
  \end{align*}
  である．

  \(Y=X\)で\(f=\id_X\)のときは，
  \begin{align*}
    \Rder{\pi_{\ast}}\muhom(G,F)
    &=\Rder{\pi_{\ast}}\muhom(G\to F)\\
    &\cong\RHOM(G,\id_{X}^{!}F)\\
    &=\RHOM(G,F)
  \end{align*}である．

  (ii) 
  命題3.4.4.\ref{prp3.4.4}より，
  \begin{align*}
    \Rder{\pi_{!}}\mu_{\diag}\RHOM(\tilde{q}_2^{-1}G,\tilde{q}_1^{!}F)
    &\cong
    \Rder{\pi_{!}}\mu_{\diag}(\rmD{G}\letens F)\\
    &\cong
    i^{-1}(\rmD{G}\letens F)\tens[] \omega_{\diag/{X\times Y}}\\
    &\cong
    i^{-1}\tilde{q}_2^{-1}\rmD{G}\tens i^{-1}\tilde{q}_1^{-1}F\tens[] \omega_{\diag/{X\times Y}}\\
    &\cong
    \varphi^{-1}\rmD{G}\tens f^{-1}F\tens[] \omega_{\diag/{X\times Y}}\\
    &\cong
    \rmD{G}\tens f^{-1}F\tens[] \omega_{\diag/{X\times Y}}
  \end{align*}
\end{proof}

関手\(\muhom\)は佐藤超局所化の関手を一般化したものである．

\begin{leftbar}
\begin{PRP}[{\cite[Prop.4.4.3]{KS90}}]\label{prp4.4.3}
  \(Y\)を\(X\)の閉部分多様体とし，
  \(j\)をうめこみ\(\conm[Y]{X}\to \cotb[X]\)とする．
  \(F\in\Dompb(X)\)とすると，次の同型が存在する．
  \[
    \muhom(A_Y,F)\cong j_{\ast}\mu_{Y}F.
  \]
\end{PRP}
\end{leftbar}
\begin{proof}
  \(f\)をうめこみ\(Y\hookrightarrow X\)とする．このとき図式\[
  \vcenter{\xymatrix{
      X\times Y
      %\ar[ld]_-{\tilde{q}_1}
      \ar@{^{(}->}[r]^-{f_2}
      \ar[d]^-{\tilde{q}_1}
      &
      X\times X
      \ar[ld]^-{q_1}
      \\
      X
      %\ar[l]^-{f}
      %\ar[u]_-{i}
      %&Y
  }}\]を考えて，
  \begin{align*}
    \mu_{\diag[X]}\RHOM(q_2^{-1}A_{Y},q_1^!F)
    &\cong
    \mu_{\diag[X]}\RG_{X\times Y}(q_1^!F)\\
    &\cong
    \mu_{\diag[X]}(\Rder{{f_2}_{\ast}}{f_2}^{!}q_1^!F)\\
    &\cong
    \mu_{\diag[X]}(\Rder{{f_2}_{\ast}}\tilde{q}_{1}^{!}F)
  \end{align*}
  である．
  \(\conm[{\diag[]}]{(X\times Y)}\)を
  \(
    Y\tpb[X]\conm[{\diag[X]}]{(X\times X)}
  \)と同一視する\footnote{
    同型\(\cotb[X]\cong\conm[{\diag[X]}]{(X\times X)}\)に引き戻し\(Y\tpb[X](\blk)\)を適用すると
    \(
      Y\tpb[X]\cotb[X]
      \cong 
      Y\tpb[X]\conm[{\diag[X]}]{(X\times X)}
    \)となる．これと\(
      Y\tpb[X]\cotb[X]
      \cong 
      \conm[{\diag[]}]{(X\times Y)}
    \)から同型を得る．
  }．座標でかくと\[    
    \vcenter{\xymatrix@C=26pt@R=10pt{
      \conm[{\diag[]}]{(X\times Y)}
      %\ar[d]_-{\tilde{f'}}
      \ar[r]^-{\sim}
      %\ar@{}[dr]|\square
      &
      Y\dpb[X]\conm[{\diag[X]}]{(X\times X)}
      %\ar[d]^-{f}
      \\
      \left(f(y),y;\xi,-f^{\ast}\xi\right)
      \ar@{}[u]|{\rotatebox{90}{\(\in\)}}
      \ar@{{|}->}[r]_-{}
      %\ar@{}[dr]|\square
      &
      \ar@{}[u]|{\rotatebox{90}{\(\in\)}}
      \left(y,f(y),f(y);\xi,-\xi\right)
    }}
  \]である．
  \begin{CMT}
    \(\mres[f_2]{\diag}\colon\diag[]\to\diag[X]\)の基本図式
    \[
      \vcenter{\xymatrix@C=36pt@R=36pt{
      T^{\ast}_{\diag[]}(X\times Y)
      \ar[dr]^-{}
      &
      \diag[]\dpb[{\diag[X]}]T^{\ast}_{\diag[X]}(X\times X)
      \ar[d]^-{}
      \ar[l]_-{\mdiff[{f_2}_{\diag[]}]}
      \ar[r]^-{\mpb[{f_2}_{\diag[]}]}
      &
      T^{\ast}_{\diag[X]}(X\times X)
      \ar[d]^-{}
      \\
      %\cotb[Y]
      &
      \diag[]
      %\ar[l]^-{}
      \ar[r]_-{\mres[{f_2}]{\diag[]}}
      &
      \diag[X]
    }}\]
    を考える．
    これは同一視\(\varphi\colon \diag[]\to Y\), 
    \(\psi\colon \diag[X]\to X\), 
    \(\conm[{\diag[X]}]{(X\times X)}\simar \cotb[X]\)のもとで
    \[
      \vcenter{\xymatrix@C=36pt@R=36pt{
      T^{\ast}_{\diag[]}(X\times Y)
      \ar[dr]^-{}
      &
      Y\dpb[X]T^{\ast}_{\diag[X]}(X\times X)
      \ar[d]^-{}
      \ar[l]_-{\mdiff[f]}
      \ar[r]^-{\mpb[f]}
      &
      \cotb[X]
      \ar[d]^-{}
      \\
      %\cotb[Y]
      &
      Y
      %\ar[l]^-{}
      \ar[r]_-{f}
      &
      X
    }}\]
    となる．
    したがって，上の同一視は\(
      \mdiff[f]=\mdiff[{f_2}_{\diag[]}]
    \)によるものである．
  \end{CMT}

  \(f_2\)は\(\diag[X]\)に対して横断的なので特に鮮明である．
  また\(f_2\)は閉うめこみなので固有的である．
  \(f_2^{-1}(\diag[X])=\diag[]\)も明らか．
  したがって命題\ref{prp4.3.4}から次が成り立つ．
  \begin{align*}
    \mu_{\diag[X]}(\Rder{{f_2}_{\ast}}\tilde{q}_{1}^{!}F)
    &\cong
    \Rder{{\mpb[{f_2}_{\diag[]}]}_{\ast}}
    {\mdiff[{f_2}_{\diag[]}]^{-1}}
    \mu_{\diag[]}(\tilde{q}_{1}^{!}F)\\
    &\cong
    \Rder{{\mpb[{f_2}_{\diag[]}]}_{\ast}}
    \mu_{\diag[]}(\tilde{q}_{1}^{!}F).
    %\quad
    &&\text{（同型\(\mdiff[{{f_2}_{\diag[]}}]\)による同一視）}
  \end{align*}
  さらに，命題\ref{prp4.3.5}
  より，
  \begin{align*}
    \Rder{{\mpb[{f_2}_{\diag[]}]}_{\ast}}\mu_{\diag[]}(\tilde{q}_{1}^{!}F)  
    &\cong
    \Rder{{\mpb[{f_2}_{\diag[]}]}_{\ast}}
    \Rder{{\mdiff[({{\widetilde{q_{1}}}_{\diag[]}})]}_{\ast}}
    \mpb[({\widetilde{q_{1}}_{\diag[]}})]^{!}
    \mu_{Y}(F)\\
    &\cong
    \Rder{{\mpb[{f_2}_{\diag[]}]}_{\ast}}
    \Rder{{\mdiff[({{\widetilde{q_{1}}}_{\diag[]}})]}_{\ast}}
    \mu_{Y}(F)
    %\quad
    &&\text{（同型\(\mpb[({\widetilde{q_{1}}_{\diag[]}})]\)による同一視）}\\
    &\cong
    j_{\ast}\mu_{Y}(F)        
    %\quad
    &&\text{（図式の可換性）}
  \end{align*}
  となる．
\end{proof}

\subsection*{ \(\muhom(G,F)\)の茎}

\(\gamma\)位相を使えば，\(\muhom(G,F)\)の茎を記述することができる．

\begin{leftbar}
\begin{PRP}[{\cite[Prop.4.4.4]{KS90}}]\label{prp4.4.4}
  \(X\)がベクトル空間であるとする．
  \(F\)と\(G\)を\(\Dompb(X)\)の対象とし，\((x_{\circ},\xi_{\circ})\in\cotb[X]\)とする．
  このとき，
  \[
    H^{j}(\muhom(G,F))_{(x_{\circ},\xi_{\circ})}
    \cong
    \indlim[U,\gamma]H^{j}(\RG(U;\RHOM(\phi_{\gamma}^{-1}\Rder{{\phi_{\gamma}}_{\ast}}G_{U},F)))
  \]
  となる．
  ただし，\(U\)は\(x_{\circ}\)の開近傍を走り，
  \(\gamma\)は\(X\)の閉凸固有錐で\(
    \gamma\subset\left\{
      v\in X;\inner<v,\xi_{\circ}><0
    \right\}\cup\{0\}
  \)をみたすもの全体を走る．
\end{PRP}
\end{leftbar}

\(\phi_{\gamma}\)は\(X\)に通常の位相を入れた空間から
\(X\)に\(\gamma\)位相を入れた空間\(X_{\gamma}\)への
自然な写像\(X\to X_{\gamma}\)を表す．

%\clearpage
\subsection*{ \(\muhom\)の関手的性質}

今度は\(\muhom\)の関手的性質をいくつか記述する．

\begin{leftbar}
\begin{PRP}[{\cite[Prop.4.4.5]{KS90}}]\label{prp.4.4.5}
  定義\ref{dfn4.4.1}の状況のとき，
  自然な射からなる以下の可換図式が得られる．
  \begin{align}
    \vcenter{\xymatrix{
      \Rder{{\mdiff[f]}_{!}}\muhom(G\to F)
      \ar[r]
      \ar[d]
      &
      \muhom(G,f^{-1}F\tens[]\omega_{Y/X})
      \ar[d]
      \\
      \Rder{{\mdiff[f]}_{\ast}}\muhom(G\to F)
      &
      \muhom(G,f^{!}F),
      \ar[l]
    }}\tag{i}\label{eq:4.4.5-1}\\
    \vcenter{\xymatrix{
      \Rder{{\mdiff[f]}_{!}}\muhom(F\leftarrow G)
      \ar[r]
      \ar[d]
      &
      \muhom(f^{!}F,G\tens[]\omega_{Y/X})
      \ar[d]
      \\
      \Rder{{\mdiff[f]}_{\ast}}\muhom(F\leftarrow G)
      &
      \muhom(f^{-1}F,G),
      \ar[l]
    }}\tag{ii}\label{eq:4.4.5-2}\\
    \vcenter{\xymatrix{
      \Rder{{\mpb[f]}_{!}}\muhom(G\to F)
      \ar[r]
      \ar[d]
      &
      \muhom(\Rder{f_{\ast}}G,F)
      \ar[d]
      \\
      \Rder{{\mpb[f]}_{\ast}}\muhom(G\to F)
      &
      \muhom(\Rder{f_{!}}G,F),
      \ar[l]
    }}\tag{iii}\label{eq:4.4.5-3}\\
    \vcenter{\xymatrix{
      \Rder{{\mpb[f]}_{!}}\muhom(F\leftarrow G)
      \ar[r]
      \ar[d]
      &
      \muhom(F,\Rder{f_{!}}G)
      \ar[d]
      \\
      \Rder{{\mpb[f]}_{\ast}}\muhom(F\leftarrow G)
      &
      \muhom(F,\Rder{f_{\ast}}G).
      \ar[l]
    }}\tag{iv}\label{eq:4.4.5-4}
  \end{align}
  \(f\)が平滑ならば，
  \eqref{eq:4.4.5-1}と\eqref{eq:4.4.5-2}の射はすべて同型となる．
  \(f\)が\(\supp{G}\)上固有ならば，
  \eqref{eq:4.4.5-3}と\eqref{eq:4.4.5-4}の射はすべて同型となる．
\end{PRP}
\end{leftbar}    

証明の前に，以降では随伴でひきおこされる射も
よく用いることに注意しておく．
例えば，命題\ref{prp4.3.5} (\ref{prp4.3.5.1}) の上辺の射
\[
  \Rder{{\mdiff[f_N]}_{!}}(
    \omega_{N/M}\tens[]{\mpb[f_N]^{-1}}\mu_{M}(F)
  )
  \to
  \mu_{N}(\omega_{Y/X}\tens[]{f^{-1}}F)     
\]
を随伴\((\Rder{{\mdiff[f_N]}_{!}},\mdiff[f_N]^{!})\)でうつした
\[
  \omega_{N/M}\tens[]{\mpb[f_N]^{-1}}\mu_{M}(F)
  \to
  \mdiff[f_N]^{!}\mu_{N}(\omega_{Y/X}\tens[]{f^{-1}}F)
\]
もよく用いる．
また，II\S6の射も頻繁に用いる．

\begin{proof}
  この節の最初の記法をここでも用いる．
  すなわち，以下の図式によって種々の射を同一視する．
  \[
  \vcenter{\xymatrix@C=36pt@R=36pt{
    \cotb[Y]
    &
    &
    Y\dpb[X]\cotb[X]
    \ar[ll]_-{\mdiff[f]}
    \\
    T^{\ast}_{\diag[Y]}(Y\times Y)
    \ar[dr]^-{}
    \ar[u]^-{\rotatebox{90}{\(\sim\)}}
    &
    \diag[Y]\dpb[{\diag[]}]T^{\ast}_{\diag[]}(X\times Y)
    \ar[d]^-{}
    \ar[l]_-{\mdiff[{f_1}_{\diag[Y]}]}
    \ar[r]^-{\mpb[{f_1}_{\diag[Y]}]}_-{\sim}
    &
    T^{\ast}_{\diag[]}(X\times Y)
    \ar[d]^-{}
    \ar[u]_-{\rotatebox{90}{\(\sim\)}}
    \\
    %\cotb[Y]
    &
    \diag[Y]
    %\ar[l]^-{}
    \ar[r]_-{\mres[f_1]{\diag[Y]}}^-{\sim}
    &
    \diag[]
  }}\]
  \[
  \vcenter{\xymatrix@C=36pt@R=36pt{
    Y\dpb[X]\cotb[X]
    \ar[rr]^-{\mpb[f]}
    &
    &
    \cotb[X]
    \\
    T^{\ast}_{\diag[]}(X\times Y)
    \ar[dr]^-{}
    \ar[u]^-{\rotatebox{90}{\(\sim\)}}
    &
    \diag[]\dpb[{\diag[X]}]T^{\ast}_{\diag[X]}(X\times X)
    \ar[d]^-{}
    \ar[l]_-{\mdiff[{f_2}_{\diag[]}]}^-{\sim}
    \ar[r]^-{\mpb[{f_2}_{\diag[]}]}
    &
    T^{\ast}_{\diag[X]}(X\times X)
    \ar[d]^-{}
    \ar[u]_-{\rotatebox{90}{\(\sim\)}}
    \\
    %\cotb[Y]
    &
    \diag[]
    %\ar[l]^-{}
    \ar[r]_-{\mres[f_2]{\diag[]}}%^-{\sim}
    &
    \diag[X]
  }}\]

  このとき，以下の自然な射が得られる．
  \begin{align}
    %(i,a)
    &\muhom(G\to F)\tag{i,a}\label{eq:4.4.5-1a}\\
    =&
    \mu_{\diag[]}
    \RHOM(\widetilde{q_2}^{-1}G,\widetilde{q_1}^{!}F)&\text{（定義）}
    \notag\\
    \to&
    \mu_{\diag[]}
    \RHOM(
      \widetilde{q_2}^{-1}G,
      \Rder{{f_1}_{\ast}}f_{1}^{-1}\widetilde{q_1}^{!}F
    )&(\id\to \Rder{{f_1}_{\ast}}f_{1}^{-1})\notag\\
    \cong&
    \mu_{\diag[]}
    \Rder{{f_1}_{\ast}}\RHOM(
      f_{1}^{-1}\widetilde{q_2}^{-1}G,
      f_{1}^{-1}\widetilde{q_1}^{!}F
    )&(f_1^{-1}\dashv \Rder{{f_1}_{\ast}})\notag\\
    \cong&
    \mu_{\diag[]}
    \Rder{{f_1}_{\ast}}\RHOM(
      {q'_2}^{-1}G,
      {q'_1}^{!}f^{-1}F
    )&\text{（図式の可換性）}\notag\\
    \cong&
    \omega_{\diag[Y]/\diag[]}
    \tens[]
    (\mpb[{f_1}_{\diag[Y]}])^{-1}\mu_{\diag[]}
    \Rder{{f_1}_{\ast}}\RHOM(
      {q'_2}^{-1}G,
      {q'_1}^{!}f^{-1}F
    )&\text{
      （\(\mres[f_1]{\diag[Y]}, \mpb[{f_1}_{\diag[Y]}]\): 同型）
    }\notag\\
    \to&
    \mdiff[{f_1}_{\diag[Y]}]^{!}\mu_{\diag[Y]}\left(
      \omega_{Y\times Y/X\times Y}
      \tens[]
      f_1^{-1}\Rder{{f_1}_{\ast}}\RHOM(
        {q'_2}^{-1}G,
        {q'_1}^{!}f^{-1}F
      )
    \right)&\text{
      （命題\ref{prp4.3.5}上辺に随伴）
    }\notag\\
    \to&
    \mdiff[{f_1}_{\diag[Y]}]^{!}\mu_{\diag[Y]}\left(
      \omega_{Y\times Y/X\times Y}
      \tens[]
      \RHOM(
        {q'_2}^{-1}G,
        {q'_1}^{!}f^{-1}F
      )
    \right)&(f_{1}^{-1}\Rder{{f_1}_{\ast}}\to \id)\notag\\
    \cong&
    \mdiff[{f_1}_{\diag[Y]}]^{!}\mu_{\diag[Y]}
    \RHOM(
      {q'_2}^{-1}G,
      {q'_1}^{!}f_{1}^{-1}F
      \tens[]
      \omega_{Y/X}    
    )\notag\\
    \cong&
    \mdiff[f]^{!}\mu_{\diag[Y]}
    \RHOM(
      {q'_2}^{-1}G,
      {q'_1}^{!}f_{1}^{-1}F
      \tens[]
      \omega_{Y/X}
    )&\text{(\(\mdiff[f]\cong\mdiff[{f_1}_{\diag[Y]}]\))}\notag
  \end{align}
  これに\(\Rder{{\mdiff[f]}_{!}}\)を施すことで上辺の射
  \begin{align*}
    \Rder{{\mdiff[f]}_{!}}\muhom(G\to F)
    &\to
    \Rder{{\mdiff[f]}_{!}}
    \mdiff[f]^{!}\micro{\diag[Y]}
    \RHOM(
      {q'_2}^{-1}G,
      {q'_1}^{!}f_{1}^{-1}F
      \tens[]
      \omega_{Y/X}
    )\\
    &\to
    \muhom(G,f_{1}^{-1}F\tens[]\omega_{Y/X})
  \end{align*}
  が得られる．
  \begin{align}
    %(i,b)
    &\muhom(G,f^{!}F)\tag{i,b}\label{eq:4.4.5-1b}\\
    =&
    \micro{\diag[Y]}
    \RHOM({q'_2}^{-1}G,{q'_1}^{!}f^{!}F)&\text{（定義）}
    \notag\\
    \cong&
    \micro{\diag[Y]}
    \RHOM(
      f_{1}^{-1}\widetilde{q_2}^{-1}G,
      f_{1}^{!}\widetilde{q_1}^{!}F
    )&\text{（図式の可換性）}
    \notag\\
    \cong&
    \micro{\diag[Y]}
    f_1^{!}\RHOM(
      \widetilde{q_2}^{-1}G,
      \widetilde{q_1}^{!}F
    )&\text{（命題3.1.13）}
    \notag\\
    \to&
    \Rder{{\mdiff[{f_1}_{\diag[Y]}]}_{\ast}}
    \mpb[{f_1}_{\diag[Y]}]^{!}\micro{\diag[]}
    \RHOM(
      \widetilde{q_2}^{-1}G,
      \widetilde{q_1}^{!}F
    )&\text{（命題\ref{prp4.3.5}下辺）}
    \notag\\
    \cong&
    \Rder{{\mdiff[f]}_{\ast}}
    \micro{\diag[]}
    \RHOM(
      \widetilde{q_2}^{-1}G,
      \widetilde{q_1}^{!}F
    )&\text{（同一視）}
    \notag
  \end{align}
  \(f\)が平滑ならば，矢印の部分は同型になる．
\end{proof}

\begin{leftbar}
\begin{PRP}[{\cite[Prop.4.4.6]{KS90}}]\label{prp.4.4.6}
    定義\ref{dfn4.4.1}の状況のとき，
    以下の自然な射が存在する．
    \begin{align}
        %(i)
        \Rder{{\mdiff[f]}_{!}}\mpb[f]^{-1}
        \muhom(\Rder{f_!}G,F)
        &\to
        \muhom(G,f^{-1}F\tens[]\omega_{Y/X}),\tag{i}\\
        %(ii)
        \Rder{{\mdiff[f]}_{!}}\mpb[f]^{-1}
        \muhom(G,\Rder{f_{\ast}}F)
        &\to
        \muhom(f^{!}F,G\tens[]\omega_{Y/X}),\tag{ii}\\
        %(iii)
        \Rder{{\mpb[f]}_{!}}\mdiff[f]^{-1}
        \muhom(G,f^{!}F)
        &\to
        \muhom(\Rder{f_{\ast}}G,F),\tag{iii}\label{eq:4.4.6-3}\\
        %(iv)
        \Rder{{\mpb[f]}_{!}}\mdiff[f]^{-1}
        \muhom(f^{-1}F,G)
        &\to
        \muhom(F,\Rder{f_{!}}G).\tag{iv}\label{eq:4.4.6-4}
    \end{align}
    \(f\)が平滑で\(\supp{G}\)上固有ならば，
    \eqref{eq:4.4.6-3}と\eqref{eq:4.4.6-4}の射は同型である．
\end{PRP}
\end{leftbar}

\(F_1\)と\(F_2\)を\(\Dompb(X)\)の対象とする．
次の標準的な射が存在する．
\begin{equation}
    \vcenter{\xymatrix{
        &\muhom(f^{!}F_2,f^{!}F_1)\ar[rd]&\\
        \muhom(f^{!}F_2,f^{-1}F_1\tens[]\omega_{Y/X})\ar[ru]\ar[rd]
        &
        &
        \muhom(f^{-1}F_2\tens[]\omega_{Y/X},f^{!}F_1)
        \\
        &\muhom(f^{-1}F_2,f^{-1}F_1)\ar[ru]&
    }}
\end{equation}
同様に\(G_1\)と\(G_2\)を\(\Dompb(Y)\)の対象とする．
次の標準的な射が存在する．
\begin{equation}
    \vcenter{\xymatrix{
        &\muhom(\Rder{f_{!}}G_2,\Rder{f_{!}}G_1)\ar[rd]&\\
        \muhom(\Rder{f_{\ast}}G_2,\Rder{f_{!}}G_1)\ar[ru]\ar[rd]
        &
        &
        \muhom(\Rder{f_{!}}G_2,\Rder{f_{\ast}}G_1)
        \\
        &\muhom(\Rder{f_{\ast}}G_2,\Rder{f_{\ast}}G_1)\ar[ru]&
    }}
\end{equation}

\begin{leftbar}
\begin{PRP}[{\cite[Prop.4.4.7]{KS90}}]\label{prp.4.4.7}
  \(F_1,F_2,G_1,G_2\)を上のようにする．
  このとき，
  自然な射で次の図式が可換になるものが存在する．
  \begin{align}
    \vcenter{\xymatrix{
      \Rder{{\mdiff[f]}_{!}}\mpb[f]^{-1}\muhom(F_2,F_1)
      \ar[r]
      \ar[d]
      &
      \muhom(f^{!}F_2,f^{-1}F_1\tens[]\omega_{Y/X})
      \ar[d]
      \\
      \Rder{{\mdiff[f]}_{\ast}}\mpb[f]^{!}(\muhom(F_2,F_1)\tens[]\omega_{Y/X}^{{\otimes}-1})
      &
      \muhom(f^{-1}F_2\tens[]\omega_{Y/X},f^{!}F_1),
      \ar[l]
    }}\tag{i}\label{eq:4.4.7-1}\\
    \vcenter{\xymatrix{
      \Rder{{\mpb[f]}_{!}}\mdiff[f]^{-1}\muhom(G_2, G_1)
      \ar[r]
      \ar[d]
      &
      \muhom(\Rder{f_{\ast}}G_2,\Rder{f_{!}}G_1)
      \ar[d]
      \\
      \Rder{{\mpb[f]}_{\ast}}\mdiff[f]^{!}(\muhom(G_2, G_1)\tens[]\omega_{Y/X})
      &
      \muhom(\Rder{f_{!}}G_2,\Rder{f_{\ast}}G_1).
      \ar[l]
    }}\tag{ii}\label{eq:4.4.7-2}
  \end{align}
  \(f\)が平滑ならば，\eqref{eq:4.4.7-1}の射はすべて同型となる．
  \(f\)が閉うめこみならば，\eqref{eq:4.4.7-2}の射はすべて同型となる．
\end{PRP}
\end{leftbar}

\subsection*{外部テンソル積}

\(\muhom\)の外部テンソル積を調べる．

\(p_X\colon X\to S\)と\(p_Y\colon Y\to S\)を多様体の射とする．
\(q_1\)と\(q_2\)で\(X\times Y\)あるいは\(X\tpb[S]Y\)上で
定義された第\(1\)射影と第\(2\)射影を表す．次を仮定する．
\begin{equation}
    \text{\(X\dpb[S]Y\)は\(X\times Y\)の部分多様体である．}
\end{equation}
\(j\)でうめこみ\(X\tpb[S]Y\hookrightarrow X\times Y\)を表す．
\((X\tpb[S]Y)\tpb[X\times Y]\cotb[(X\times Y)]\cong\cotb[X]\tpb[S]\cotb[Y]\)より，次の\(2\)つの射が得られる．
\begin{equation}
    \cotb[{\left(X\dpb[S]Y\right)}]
    \underset{{\mdiff[j]}}{\longleftarrow}
    \cotb[X]\dpb[S]\cotb[Y]
    \underset{{\mpb[j]}}{\longrightarrow}
    \cotb[X]\times\cotb[Y].
\end{equation}

\begin{leftbar}
\begin{PRP}[{\cite[Prop.4.4.8]{KS90}}]
  \(F_1\)と\(F_2\)を\(\Dompb(X)\)の対象とし
  \(G_1\)と\(G_2\)を\(\Dompb(Y)\)の対象とする．
  次の標準的な射が存在する．
  \begin{align}
    \tag{i}
    \Rder{{\mdiff[j]}_{!}}\left(
      \muhom(F_2,F_1)\lletens[S]\muhom(G_2,G_1)
    \right)&\\
    \to
    \muhom&\left(
      F_2\lletens[S]G_2,F_1\lletens[S]G_1
    \right),\notag\\
    \tag{ii}
    \Rder{{\mdiff[j]}_{!}}\left(
      \muhom(F_2,F_1)^{a}\lletens[S]\muhom(G_2,G_1)
    \right)&\\
    \to
    \muhom&(
      \RHOM(q_{1}^{-1}F_1,q_{2}^{-1}G_2),
      \RHOM(q_{1}^{-1}F_2,q_{2}^{-1}G_1)
    )\notag.
  \end{align}
\end{PRP}
\end{leftbar}

\subsection*{\(\muhom\)の合成}
関手\(\muhom\)を合成する．

\(g\colon Z\to Y\)と\(f\colon Y\to X\)を多様体の射とし，
\(h=f\circ g\)とする．
\(q_i\), \(q'_i\), \(q''_i\) \((i=1,2)\)をそれぞれ
\(Z\times Y\), \(X\times Z\), \(Y\times Z\)の第\(i\)射影とし，
\(q_{ij}\)を\(X\times Y\times Z\)の第\((i,j)\)射影とする．
\(\widetilde{q}_{ij}\)で\(X\times Y\times Y\times Z\)の
第\((i,j)\)射影を表し，
\(p_{ij}\)と\(\widetilde{p}_{ij}\)でそれぞれ
\(\cotb[X]\times\cotb[Y]\times\cotb[Z]\)と
\(\cotb[X]\times\cotb[Y]\times\cotb[Y]\times\cotb[Z]\)の
第\((i,j)\)射影を表す．
\(r_{ij}\)で\(\widetilde{p}_{ij}\)の\(
  \diag[Y]\tpb[Y\times Y]\cotb[(X\times Y\times Y\times Z)]
\)への制限を表す．
最後に\(j\)で対角うめこみ\(
  X\times Y\times Z
  \hookrightarrow 
  X\times Y\times Y\times Z
\)を表す．\(\mdiff[j]\)は同型
\begin{equation}
  \diag[Y]\dpb[Y\times Y]
  \conm[{\diag[f]\times\diag[g]}]{(X\times Y\times Y\times Z)}
  \overunderset{\sim}{{\mdiff[j]}}{\longrightarrow}
  \conm[{\diag[f]\tpb[Y]\diag[g]}]{(X\times Y\times Z)}
\end{equation}
をひきおこし，
\(\mpb[{q_{13}}]\)は同型
\begin{equation}
  \left(
    \diag[f]
    \dpb[Y]
    \diag[g]
  \right)\dpb[{\diag[h]}]
  \conm[{\diag[h]}]{(X\times Z)}
  \overunderset{\sim}{{\mpb[{q_{13}}]}}{\longrightarrow}
  \conm[{\diag[h]}]{(X\times Z)}
\end{equation}
をひきおこす．
これによって次の可換図式を得る．

\begin{equation}
  \vcenter{\xymatrix@C=36pt@R=36pt{
    &
    &
    \conm[{\diag[f]}]{(X\times Y)}
    \cong
    Y\dpb[X]\cotb[X]
    \\
    \diag[Y]\dpb[Y\times Y]
    \conm[{\diag[f]\times\diag[g]}]{(X\times Y\times Y\times Z)}
    \ar[urr]^-{r_{12}}
    \ar[drr]_-{r_{34}}
    \ar[r]_-{\mpb[j]}^-{\sim}
    &
    \conm[{\diag[f]\tpb[Y]\diag[g]}]{(X\times Y\times Z)}
    &
    \conm[{\diag[h]}]{(X\times Z)}
    \cong
    Z\dpb[X]\cotb[X]
    \ar@{_{(}->}[l]^-{\mdiff[{q_{13}}]}
    %T^{\ast}_{\diag[X]}(X\times X)
    \ar[d]^-{\mdiff[f]}
    \ar[u]_-{\mpb[g]}
    \\
    %\cotb[Y]
    &
    %\diag[]
    %\ar[l]^-{}
    %\ar[r]_-{\mres[f_2]{\diag[]}}%^-{\sim}
    &
    \conm[{\diag[g]}]{(Y\times Z)}
    \cong
    Z\dpb[Y]\cotb[Y].
  }}
\end{equation}

\begin{CMT}
  ここで\(\mpb[g]\)と書いてある射は\(
    g\colon Z\to Y
  \)に\((\blk)\tpb[X]\cotb[X]\)を適用したものである．
  \(\mdiff[f]\)の方は，\(f\colon Y\to X\)に\(Z\tpb[Y](\blk)
  \)を適用したものである．
\end{CMT}

\begin{leftbar}
\begin{PRP}[{\cite[Prop.4.4.9]{KS90}}]
  \(F\in\Dompb(X)\), \(G\in\Dompb(Y)\), 
  \(H\in\Dompb(Z)\)とする．
  次の標準的な射が定まる．
  \[
    \mdiff[f]^{-1}\muhom(H\to G)
    \lltens[]
    \mpb[g]^{-1}\muhom(G\to F)
    \to
    \muhom(H\to F).
  \]
\end{PRP}
\end{leftbar}

\begin{proof}
  次のように記号を定める．
  \begin{align*}
    K_1&\coloneqq \RHOM(q_{2}^{-1}G,q_{1}^{!}F),\\
    K_2&\coloneqq \RHOM({q_{2}''}^{-1}H,{q_{1}''}^{!}G),\\
    L&\coloneqq 
    \mpb[g]^{-1}\micro{\diag[f]}(K_1)
    \lltens[]
    \mdiff[f]^{-1}\micro{\diag[g]}(K_2)
  \end{align*}
  このとき，
  \begin{align*}
    L&\simar 
    \mdiff[{q_{13}}]^{-1}{\mdiff[j]}_{\ast}r_{12}^{-1}\micro{\diag[f]}(K_1)
    \lltens[]
    \mdiff[{q_{13}}]^{-1}{\mdiff[j]}_{\ast}r_{34}^{-1}\micro{\diag[g]}(K_2)\\
    &\simar 
    \mdiff[{q_{13}}]^{-1}{\mdiff[j]}_{\ast}\left(r_{12}^{-1}\micro{\diag[f]}(K_1)
    \lltens[]
    r_{34}^{-1}\micro{\diag[g]}(K_2)\right)&&\text{（(2.6.22)が同型．(2.6.18)）}\\
    &\cong
    \mdiff[{q_{13}}]^{-1}{\mdiff[j]}_{\ast}
    \left(
      \mpb[j]^{-1}\tilde{p}_{12}^{-1}\micro{\diag[f]}(K_1)
      \lltens[]
      \mpb[j]^{-1}\tilde{p}_{34}^{-1}\micro{\diag[g]}(K_2)
    \right)
    &&\text{（図式の可換性）}\\
    &\simar
    \mdiff[{q_{13}}]^{-1}{\mdiff[j]}_{\ast}\mpb[j]^{-1}
    \left(
      \tilde{p}_{12}^{-1}\micro{\diag[f]}(K_1)
      \lltens[]
      \tilde{p}_{34}^{-1}\micro{\diag[g]}(K_2)
    \right)
    &&\text{（(2.6.18)）}\\
    &\to
    \mdiff[{q_{13}}]^{-1}{\mdiff[j]}_{\ast}\mpb[j]^{-1}
    \micro{\diag[f]\times\diag[g]}\left(
      \tilde{q}_{12}^{-1}K_1\lltens[]\tilde{q}_{34}^{-1}K_2
    \right)
    &&\text{（命題\ref{prp4.3.6}）}\\
    &\to
    \mdiff[{q_{13}}]^{-1}
    \micro{\diag[f]\dpb[Y]\diag[g]}\left(
      j^{-1}\left(
        \tilde{q}_{12}^{-1}K_1\lltens[]\tilde{q}_{34}^{-1}K_2
      \right)
    \right)
    &&\text{（命題\ref{prp4.3.5} \eqref{prp4.3.5.3}）}\\
    &\simar
    \mdiff[{q_{13}}]^{-1}
    \micro{\diag[f]\dpb[Y]\diag[g]}\left(
      j^{-1}\tilde{q}_{12}^{-1}K_1
      \lltens[]
      j^{-1}\tilde{q}_{34}^{-1}K_2
    \right)
    &&\text{（(2.6.18)）}\\
    &\simar
    \mdiff[{q_{13}}]^{-1}
    \micro{\diag[f]\dpb[Y]\diag[g]}\left(
      q_{12}^{-1}K_1
      \lltens[]
      q_{23}^{-1}K_2
    \right)
    &&\text{（図式の可換性）}\\
    &\to
    \mdiff[{q_{13}}]^{-1}
    \micro{\diag[f]\dpb[Y]\diag[g]}\left(      
      \RHOM(
        q_{13}^{-1}{q_{2}'}^{-1}H,
        q_{13}^{!}{q'_{1}}^{!}F)
    \right)
    &&\text{（後述）}\tag{\(\natural\)}\label{4.4.9.hom}\\
    &\cong
    \Rder{{\mdiff[{q_{13}}]}_{!}}\mdiff[{q_{13}}]^{-1}
    \micro{\diag[f]\dpb[Y]\diag[g]}\left(      
      \RHOM(
        q_{13}^{-1}{q_{2}'}^{-1}H,
        q_{13}^{!}{q'_{1}}^{!}F)
    \right)
    &&\text{（同一視から）}\\
    &\to
    \micro{\diag[f]\dpb[Y]\diag[g]}\left(
      \Rder{{q_{13}}_{!}}      
      \RHOM(
        q_{13}^{-1}{q_{2}'}^{-1}H,
        q_{13}^{!}{q'_{1}}^{!}F)
    \right)
    &&\text{（命題\ref{prp4.3.4}上辺）}\\
    &\to
    \micro{\diag[f]\dpb[Y]\diag[g]}\left(
      \RHOM(
        \Rder{{q_{13}}_{!}}q_{13}^{-1}{q_{2}'}^{-1}H,
        \Rder{{q_{13}}_{!}}q_{13}^{!}{q'_{1}}^{!}F)
    \right)
    &&\text{（(2.6.26)）}\\
    &\to
    \micro{\diag[f]\dpb[Y]\diag[g]}\left(
      \RHOM({q_{2}'}^{-1}H,{q'_{1}}^{!}F)
    \right)
    &&\text{（随伴の単位・余単位）}\\
    &=
    \muhom(H\to F)
  \end{align*}
  となり，主張が示される．

  \eqref{4.4.9.hom}の部分について説明する．
  \(
    \omega_{1}
    \coloneqq
    \omega_{X\times Y\times Z/X\times Y}
  \)とおく．
  自然な射\[
    q_{12}^{-1}K_1
    \lltens[]
    \omega_{1}
    \to q_{12}^{!}K_1
  \]の両辺に\(
    \omega_{1}^{\otimes-1}
  \)をかけると射
  \(
    q_{12}^{-1}K_1
    \to q_{12}^{!}K_1
    \lltens[]
    \omega_{1}^{\otimes-1}
  \)が定まる．次に
  \(\omega_{2}\coloneqq \omega_{Y\times Z/Y}\)
  とおく．このとき，
  \(
    q_{23}^{-1}\omega_2\lltens[]\omega_1
    \cong 
    A_{X\times Y\times Z}
  \)が成り立つ．
  最後に
  \begin{align*}
    \RHOM(F,G)\lltens[]F&\to G, 
    \tag{{\cite[(2.2.10)]{KS90}}}\\
    \RHOM(F,G)\lltens[]H&\to \RHOM(F,G\lltens[]H), 
    \tag{{\cite[(2.2.11)]{KS90}}}
  \end{align*}
  から，自然な射
  \begin{align*}
    %\vcenter{
      \RHOM(F,G)\lltens[]\RHOM(G,H)
      &\to \RHOM(F,G\lltens[]\RHOM(G,H)) \\
      &\to \RHOM(F,H)
    %}    
    \tag{\(\natural\natural\)}\label{4.4.9.contraction}
  \end{align*}
  が定まる．以上の3つの準備のもとで，
  以下の射が得られる．
  \begin{align*}
    &q_{12}^{-1}K_1
      \lltens[]
      q_{23}^{-1}K_2
    &&\text{}\\
    \to&
      q_{12}^{!}K_1
      \lltens[]\omega_{1}^{\otimes-1}
      \lltens[]
      q_{23}^{-1}K_2
    &&\text{}\\
    \to&
      \RHOM(q_{12}^{-1}q_{2}^{-1}G,q_{12}^{!}q_{1}^{!}F)
      \lltens[]\omega_{1}^{\otimes-1}
      \lltens[]
      \RHOM(q_{23}^{-1}{q_{2}''}^{-1}H,q_{23}^{-1}{q_{1}''}^{!}G)
    &&%\text{（(2.6.27)と命題3.1.13）}
    \\
    \to&
      \RHOM(q_{12}^{-1}q_{2}^{-1}G,q_{13}^{!}{q'_{1}}^{!}F)
      \lltens[]\omega_{1}^{\otimes-1}
      \lltens[]
      \RHOM(q_{13}^{-1}{q_{2}'}^{-1}H,q_{23}^{-1}({q_{1}''}^{-1}G\lltens[]\omega_{2}))
    &&\text{}\\
    \simar&
      \RHOM(q_{12}^{-1}q_{2}^{-1}G,q_{13}^{!}{q'_{1}}^{!}F)
      \lltens[]
      \omega_{1}^{\otimes-1}
      \lltens[]
      \RHOM(q_{13}^{-1}{q_{2}'}^{-1}H,q_{23}^{-1}{q_{1}''}^{-1}G\lltens[]q_{23}^{-1}\omega_{2})
    &&\text{}\\
    \simar&
      \RHOM(q_{12}^{-1}q_{2}^{-1}G,q_{13}^{!}{q'_{1}}^{!}F)
      \lltens[]
      \RHOM(
        q_{13}^{-1}{q_{2}'}^{-1}H,
        q_{12}^{-1}q_{2}^{-1}G
      )
    \lltens[]
    \omega_{1}^{\otimes-1}
    \lltens[]q_{23}^{-1}\omega_{2}
    &&\text{}\\
    \cong&
      \RHOM(q_{12}^{-1}q_{2}^{-1}G,q_{13}^{!}{q'_{1}}^{!}F)
      \lltens[]
      \RHOM(
        q_{13}^{-1}{q_{2}'}^{-1}H,
        q_{12}^{-1}q_{2}^{-1}G
      )
    &&\text{}\\
    \underset{\eqref{4.4.9.contraction}}{\to}&
      \RHOM(
        q_{13}^{-1}{q_{2}'}^{-1}H,
        q_{13}^{!}{q'_{1}}^{!}F)
    &&\text{}
  \end{align*}
  したがって，以上の射の合成に関手\(
    \mdiff[{q_{13}}]^{-1}
    \micro{\diag[f]\dpb[Y]\diag[g]}(\blk)
  \)を適用することで\eqref{4.4.9.hom}が定まる．
\end{proof}

\begin{leftbar}
\begin{CRL}[{\cite[Cor.4.4.10]{KS90}}]
  \(F_1,F_2,F_3\in\Dompb(X)\)とする．
  このとき次の標準的な射が定まる．
  \[
    \muhom(F_1, F_2)
    \lltens[]
    \muhom(F_2, F_3)
    \to
    \muhom(F_1, F_3).
  \]
\end{CRL}
\end{leftbar}

\begin{leftbar}
\begin{PRP}[{\cite[Prop.4.4.11]{KS90}}]
  \(K_1, F_1\in\Dompb(X\times Y)\), 
  \(F_2,K_2\in\Dompb(Y\times Z)\)
  に対し，次の標準的な射が定まる．
  \begin{equation}
    \Rder{{p_{13^{a}}}_{!}}
    \left(
      p_{12^{a}}^{-1}\muhom(K_1,F_1)
      \lltens[]
      p_{23^{a}}^{-1}\muhom(K_2,F_2)
    \right)
    \to
    \muhom(K_1\circ K_2,F_1\circ F_2),
  \end{equation}
  ただし，\(K_1\circ K_2,F_1\circ F_2\)は核の合成である．
\end{PRP}
\end{leftbar}
















\clearpage
% ======================================================
% appendix
% ======================================================
\appendix
\section{ノート中で用いる主張}
\begin{leftbar}
\begin{PRP}[{\cite[Prop.3.4.4.]{KS90}}]\label{prp3.4.4}
  \(X\)と\(Y\)を有限c柔軟次元局所コンパクト（ハウスドルフ）空間とし，
  \(q_1\)と\(q_2\)をそれぞれ\(X\times Y\)から\(X\)と\(Y\)への射影とする．
  \(F\in\Dompb(A_X)\), \(G\in\Dompb(A_Y)\)とし，\(F\)はコホモロジー構成可能層とする．
  このとき，
  \[
    \vD{F}\letens G\to \RHOM(q_1^{-1}F,q_2^{!}G)
  \]は同型である．
  \(X\)が\(C^0\)多様体ならば，\[
    \vD'{F}\letens G\to \RHOM(q_1^{-1}F,q_2^{-1}G)
  \]も同型である．
\end{PRP}
\end{leftbar}

\begin{leftbar}
\begin{PRP}[{\cite[Prop.4.3.4]{KS90}}]\label{prp4.3.4}
  \(G\in\Dompb(Y)\)とする．
  このとき，標準的な射による可換図式    
  \[    
      \vcenter{\xymatrix@C=36pt@R=36pt{
      \Rder{{\mpb[f_{N}]}_{!}}{\mdiff[f_{N}]^{-1}}\mu_{N}(G)
      \ar[d]
      \ar[r]
      %\ar@{}[dr]|\square
      &
      \mu_{M}(\Rder{f_{!}}G)
      %\ar[r]
      \ar[d]
      \\
      \Rder{{\mpb[f_{N}]}_{\ast}}({\mdiff[f_{N}]^{!}}\mu_{N}(G)\tens[]\omega_{Y/X}\tens[]\omega_{N/M}^{\otimes{-1}})
      %\ar@{}[dr]|\square
      &
      \mu_{M}(\Rder{f_{\ast}}G)
      \ar[l]
      }}
  \]が存在する．
  さらに，\(\supp(G)\to X\)と\(
      C_N(\supp(G))\to T_{M}X
  \)が固有かつ\(\supp(G)\cap f^{-1}(M)\subset N\)ならば，
  以上の射はすべて同型である．
  
  とくに\(f^{-1}(M)=N\)かつ，
  \(f\)が\(M\)に関して鮮明\footnotemark[2] (clean) かつ\(\supp(G)\)上固有のとき，
  以上の射はすべて同型である．
\end{PRP}
\end{leftbar}

\footnotetext[2]{試験的な訳語}

\begin{leftbar}
\begin{PRP}[{\cite[Prop.4.3.5]{KS90}}]\label{prp4.3.5}
  \(F\in\Dompb(X)\)とする．
  \begin{enumerate}[(i)]
      \item 標準的な射で次の図式を可換にするものが存在する．\label{prp4.3.5.1}
      \[\vcenter{\xymatrix@C=56pt@R=46pt{
          \Rder{{\mdiff[f_N]}_{!}}(\omega_{N/M}\tens[]{\mpb[f_N]^{-1}}\mu_{M}(F))
          \ar[d]^{}
          \ar[r]_{}
          %\ar@{}[dr]|\square
          &
          \mu_{N}(\omega_{Y/X}\tens[]{f^{-1}}F)
          %\ar[r]
          \ar[d]
          \\
          \Rder{{\mdiff[f_N]}_{\ast}}{\mpb[f_N]^{!}}\mu_{M}(F)
          %\ar@{}[dr]|\square
          &
          \mu_{N}({f^{!}}F)
          \ar[l]_{\beta}
      }}\]
      ここで縦向きの矢印は(3.1.6)から導かれるものである．
      \item \label{prp4.3.5.2}
      \(f\colon Y\to X\)と\(\mres[f]{N}\colon N\to M\)が
      平滑ならば，これらの射はすべて同型となる．    
      \item \label{prp4.3.5.3}
      \(f\)が\(M\)に横断的かつ\(f^{-1}(M)=N\)ならば，
      自然な射\[
          \Rder{{\mdiff[f_N]}_{\ast}}\mpb[f_N]^{-1}\mu_{M}(F)
          \to \mu_{N}(f^{-1}F)
      \]が存在する．
  \end{enumerate}
\end{PRP}
\end{leftbar}

\begin{leftbar}
\begin{PRP}[{\cite[Prop.4.3.6.]{KS90}}]\label{prp4.3.6}
  外部テンソル積
\end{PRP}
\end{leftbar}














%===============================================
% 参考文献スペース
%===============================================
\begin{thebibliography}{20} 
  \bibitem[dR]{dR} 
  de Rham, \textit{Vari\'et\'es diff\'erentiables}, Hermann, 1953.
  和訳：ド・ラーム, 微分多様体, 東京図書, 1974.
  \bibitem[GKS12]{GKS12} St\'ephane Guillermou, 
  Masaki Kashiwara, Pierre Schapira, 
  \textit{Sheaf Quantization of Hamiltonian Isotopies and Applications to Nondisplaceability Problems}, 
  Duke. Math. J. 161, No. 2, pp.201--245, 2012.  
  \bibitem[GP74]{GP74} Victor Guillemin, Alan Pollack, 
    \textit{Differential Topology}, 
    Prentice-Hall, 1974.
  \bibitem[KS90]{KS90} Masaki Kashiwara, Pierre Schapira, 
    \textit{Sheaves on Manifolds}, 
    Grundlehren der Mathematischen Wissenschaften, 292, Springer, 1990.
  \bibitem[Hat89]{Hat89} 服部晶夫, 多様体 増補版, 岩波全書 288, 岩波書店, 1989.
  %\bibitem[Og02]{Og02} 小木曽啓示, 代数曲線論, 朝倉書店, 2022.
  \bibitem[Hor63]{Hor63} Lars H\"ormander, 
  \textit{Linear Pertial Differential Operators}, Fourth Printing, 
  Grundlehren der Mathematischen Wissenschaften, 116, Springer, 1963.
  \bibitem[Hor71]{Hor71} Lars H\"ormander, 
  \textit{Fourier Integral Operators I}, 
  Acta Math. 127, 79--183, (1971).
  \bibitem[Hor89]{Hor89} Lars H\"ormander, 
  \textit{The Analysis of Linear Pertial Differential Operators I}, Second Edition,
  Grundlehren der Mathematischen Wissenschaften, 256, Springer, 1989.
  \bibitem[Hor85]{Hor85} Lars H\"ormander, 
  \textit{The Analysis of Linear Pertial Differential Operators III}, 
  Grundlehren der Mathematischen Wissenschaften, 274, Springer, 1985.
  \bibitem[Sch85]{Sch85} Pierre Schapira, 
    \textit{Microdifferential Systems in the Complex Domain}, 
    Grundlehren der Mathematischen Wissenschaften, 269, 
    Springer, 1985.
    \bibitem[ST92]{ST92} Pierre Schapira, Nobuyuki Tose, 
    \textit{Morse Inequalities for \(\rr\)-constructible Sheaves}, 
    Adv. Math. 93, pp.1--8, 1992. 
    \bibitem[Zwo12]{Zwo12} Maciej Zworski,
    \textit{Semiclsaaical Analysis}, 
    Graduate Studies in Mathematics, 138, 
    American Mathematical Society, 
    2012. 
\end{thebibliography}

%===============================================




\end{document}
